\chapter{1955年全国卷}

\section{解答题}

\FigureLayoutDeclare{img/1955/national/q1_solution_fig1.png}{2.8cm}{0bp 0bp 0bp 0bp}
\FigureLayoutDeclare{img/1955/national/q1_solution_fig2.png}{4.4cm}{0bp 0bp 0bp 0bp}
\FigureLayoutDeclare{img/1955/national/q3_solution_fig1.png}{5.8cm}{0bp 0bp 0bp 0bp}

\begin{problem}
\begin{enumerate}
\item 以二次方程 \(x^{2}-3x-1=0\) 的两根的平方为两根作一二次方程.
\item 等腰三角形一腰的长是底边的 \(4\) 倍，求这三角形各角的余弦.
\item 已知正四棱锥底边的长为 \(a\)，侧棱与底面的交角为 \(45^{\circ}\)，求这棱锥的高.
\item 写出：二面角的平面角的定义.
\end{enumerate}
\end{problem}

\begin{solution}
\begin{enumerate}
\item
设方程 \(x^{2}-3x-1=0\) 的两根为 \(\alpha\)、\(\beta\)，由韦达定理得 \(\alpha+\beta=3\)，\(\alpha\beta=-1\). 以 \(\alpha^{2}\)、\(\beta^{2}\) 为根的二次方程的根和为 \(\alpha^{2}+\beta^{2}=\left(\alpha+\beta\right)^{2}-2\alpha\beta=3^{2}-2\left(-1\right)=11\)，根积为 \(\alpha^{2}\beta^{2}=\left(\alpha\beta\right)^{2}=1\). 因此所求方程为 \(y^{2}-11y+1=0\).

\item
设 \(\triangle ABC\) 中 \(AB=AC=4BC\)，\(AD\perp BC\). 由于三角形为等腰三角形，\(D\) 是 \(BC\) 的中点，因此 \(BD=DC=\frac{1}{2}BC\).

\solutionfigure{\bitmapfigure[max width=2.4cm]{img/1955/national/q1_solution_fig1.png}}

由余弦定理，\(\cos A=\dfrac{AB^{2}+AC^{2}-BC^{2}}{2AB\cdot AC}=\dfrac{16BC^{2}+16BC^{2}-BC^{2}}{2\cdot 4BC\cdot 4BC}=\dfrac{31}{32}\). 在直角三角形 \(ACD\) 中，\(\cos C=\dfrac{CD}{CA}=\dfrac{\frac{1}{2}BC}{4BC}=\dfrac{1}{8}\)，又因 \(B=C\)，所以 \(\cos B=\cos C=\dfrac{1}{8}\).

\item
设 \(S-ABCD\) 为正四棱锥，\(H\) 为正方形底面 \(ABCD\) 的中心，\(SH\) 为棱锥的高. 由于 \(SH\perp\) 底面，\(AH\) 是侧棱 \(SA\) 在底面上的射影，所以侧棱与底面的交角为 \(\angle SAH=45^{\circ}\).

\solutionfigure{\bitmapfigure[max width=3.6cm]{img/1955/national/q1_solution_fig2.png}}

在直角三角形 \(SHA\) 中，\(\angle SHA=90^{\circ}\)，且 \(\angle SAH=45^{\circ}\)，故 \(\angle ASH=45^{\circ}\)，从而 \(\triangle SHA\) 为等腰直角三角形，\(SH=AH\). 正方形底面的对角线为 \(a\sqrt{2}\)，而 \(H\) 是其中心，所以 \(AH=\dfrac{a\sqrt{2}}{2}\). 因此棱锥的高为 \(SH=\dfrac{a\sqrt{2}}{2}\).

\item
在二面角的棱上任取一点，在二面角的两个面内分别过该点作棱的垂线；这两条垂线在两个面内所夹的角，叫做这个二面角的平面角.
\end{enumerate}
\end{solution}

\begin{problem}
求 \(b\)，\(c\)，\(d\) 的值，使多项式 \(x^{3}+bx^{2}+cx+d\) 适合下列三条件：

\begin{circlelist}
\item 被 \(x-1\) 整除；
\item 被 \(x-3\) 除时余 \(2\)；
\item 被 \(x+2\) 除与被 \(x-2\) 除时余数相等.
\end{circlelist}
\end{problem}

\begin{solution}
设 \(P(x)=x^{3}+bx^{2}+cx+d\). 由余式定理，\(P(x)\) 被 \(x-1\) 整除，所以 \(P(1)=0\)，即 \(1+b+c+d=0\). 又因 \(P(x)\) 被 \(x-3\) 除时余数为 \(2\)，所以 \(P(3)=2\)，即 \(27+9b+3c+d=2\). 第三个条件给出 \(P(-2)=P(2)\)，即 \(-8+4b-2c+d=8+4b+2c+d\).

化简第三个等式得 \(-16-4c=0\)，所以 \(c=-4\). 代入前两个等式，分别得到 \(b+d=3\) 和 \(9b+d=-13\). 两式相减得 \(8b=-16\)，从而 \(b=-2\)，再由 \(b+d=3\) 得 \(d=5\). 因此 \((b,c,d)=(-2,-4,5)\).
\end{solution}

\begin{problem}
由直角三角形勾上一点 \(D\) 作弦 \(AB\) 的垂线交弦于 \(E\)，股的延长线于 \(F\)，外接圆周于 \(Q\)，求证：\(EQ\) 为 \(EA\) 与 \(EB\) 的比例中项，又为 \(ED\) 与 \(EF\) 的比例中项.
\end{problem}

\begin{solution}
\solutionfigure{\bitmapfigure[max width=4.6cm]{img/1955/national/q3_solution_fig1.png}}
由题意，\(\triangle ABC\) 在 \(C\) 处为直角，\(AB\) 是其斜边，\(D\) 在直角边 \(AC\) 上，且 \(DE\perp AB\)，\(F\) 在直角边 \(BC\) 的延长线上，\(Q\) 是直线 \(DE\) 与外接圆的交点. 因为 \(\triangle ABC\) 是直角三角形，\(AB\) 是其外接圆的直径，所以 \(\angle AQB=90^{\circ}\). 又 \(QE\perp AB\)，故 \(QE\) 是直角三角形 \(AQB\) 斜边 \(AB\) 上的高.

连结 \(QA\)、\(QB\). 由 \(\triangle AQE\sim\triangle QEB\)，得 \(\dfrac{EA}{QE}=\dfrac{QE}{EB}\)，所以 \(EQ^{2}=EA\cdot EB\). 这说明 \(EQ\) 是 \(EA\) 与 \(EB\) 的比例中项.

再证第二个结论. 因为 \(AD\) 与 \(AC\) 共线，所以 \(\angle EAD=\angle CAB\). 又 \(BF\) 与 \(BC\) 共线，故 \(\angle ABF=\angle ABC\). 在直角三角形 \(ABC\) 中，\(\angle CAB+\angle ABC=90^{\circ}\). 另一方面，\(FE\perp AB\)，所以 \(\angle EFB+\angle ABF=90^{\circ}\). 因而 \(\angle EAD=\angle EFB\)，并且 \(\angle AED=\angle FEB=90^{\circ}\)，从而 \(\triangle AED\sim\triangle FEB\).

对应边成比例，故 \(\dfrac{EA}{EF}=\dfrac{ED}{EB}\)，即 \(EA\cdot EB=ED\cdot EF\). 结合前面得到的 \(EQ^{2}=EA\cdot EB\)，可得 \(EQ^{2}=ED\cdot EF\). 由于各线段长度均为正，\(EQ\) 是 \(ED\) 与 \(EF\) 的比例中项.
\end{solution}

\begin{problem}
解方程 \(\cos 2x=\cos x+\sin x\)，求 \(x\) 的通值.
\end{problem}

\begin{solution}
由 \(\cos 2x=\cos^{2}x-\sin^{2}x\)，原方程等价于 \(\cos^{2}x-\sin^{2}x-\cos x-\sin x=0\). 因式分解得 \(\left(\cos x+\sin x\right)\left(\cos x-\sin x-1\right)=0\).

因此分两种情况讨论.

当 \(\cos x+\sin x=0\) 时，\(\sqrt{2}\sin\left(x+\dfrac{\pi}{4}\right)=0\)，所以 \(x+\dfrac{\pi}{4}=n\pi\)，即 \(x=n\pi-\dfrac{\pi}{4}\)，其中 \(n\in\Z\).

当 \(\cos x-\sin x-1=0\) 时，\(\cos x-\sin x=1\). 两边同除以 \(\sqrt{2}\)，得 \(\cos\left(x+\dfrac{\pi}{4}\right)=\dfrac{1}{\sqrt{2}}=\cos\dfrac{\pi}{4}\).

所以 \(x+\dfrac{\pi}{4}=2n\pi\pm\dfrac{\pi}{4}\)，从而 \(x=2n\pi\) 或 \(x=2n\pi-\dfrac{\pi}{2}\)，其中 \(n\in\Z\).

综上，方程的通解为 \(x=n\pi-\dfrac{\pi}{4}\)，或 \(x=2n\pi\)，或 \(x=2n\pi-\dfrac{\pi}{2}\)，其中 \(n\in\Z\).
\end{solution}

\begin{problem}
一三角形三边的长成等差级数，其周长为 \(12\text{ 尺}\)，面积为 \(6\text{ 平方尺}\)，求证这三角形为一直角三角形.
\end{problem}

\begin{solution}
设三角形三边按等差数列的顺序为 \(x-y\)、\(x\)、\(x+y\)（单位均为 \(\text{尺}\)），其中 \(x\) 是中项，\(y\) 是公差. 由周长为 \(12\text{ 尺}\)，得 \((x-y)+x+(x+y)=12\)，所以 \(3x=12\)，即 \(x=4\).

三角形半周长为 \(s=6\text{ 尺}\). 由海伦公式 \(S^{2}=s\left(s-(x-y)\right)\left(s-x\right)\left(s-(x+y)\right)\)，代入面积 \(S=6\text{ 平方尺}\) 和 \(s=6\text{ 尺}\)，得 \(36=6(6-x+y)(6-x)(6-x-y)\).

代入 \(x=4\)，得 \(36=6(2+y)\cdot 2\cdot(2-y)=12(4-y^{2})\)，所以 \(y^{2}=1\)，即 \(y=\pm1\). 当 \(y=1\) 时三边为 \(3\text{ 尺}\)、\(4\text{ 尺}\)、\(5\text{ 尺}\)；当 \(y=-1\) 时三边只是顺序相反，仍为 \(3\text{ 尺}\)、\(4\text{ 尺}\)、\(5\text{ 尺}\).

最后，\(3^{2}+4^{2}=9+16=25=5^{2}\). 由勾股定理的逆定理，最长边为 \(5\text{ 尺}\) 的这个三角形为直角三角形.
\end{solution}
